\SourceLeaf{324}%
oder wenn man für \(u\) den Wert aus (3) einsetzt und nach (8), (9) und (11)
\[
s=\frac{r\cos\theta^2}{a\cos\varphi}=\frac{r\cos\theta\operatorname{tang}P}{a\operatorname{tang}Q}
\]
setzt, erhält man
\[
t=\frac{r\cos\theta\operatorname{tang}P\mathord{.}\sin U}{a}(1+\tfrac1{24}uu+\tfrac1{12}ss-\frac{ee\cos P^2}{6\cos\varphi^2}\mathord{.}qq)
\]
Für \(\cos\theta\mathord{.}\operatorname{tang}P=\BareRoot{(1-ee\sin P^2)}\mathord{.}\operatorname{tang}P\) erhält man nach (4) den bis zur hier benötigten Ordnung entwickelten Wert
\[
\BareRoot{(1-ee\sin B^2)}\mathord{.}\operatorname{tang}B(1+\frac{3ee-6e^4\sin P^2+3e^4\sin P^4}{8\cos\varphi^2(1-ee\sin P^2)}\mathord{.}qq)
\]
und folglich
\[
\begin{aligned}
t={}&\frac{rk\operatorname{tang}B\sin U}{a}(1+\frac{5ee+(4ee-14e^4)\sin P^2+5e^4\sin P^4}{24\cos\varphi^2(1-ee\sin P^2)}\mathord{.}qq\\
&\qquad+\tfrac1{12}ss+\tfrac1{24}uu)
\end{aligned}
\]
Nimmt man hier nun noch dieselben Ersetzungen wie im vorigen Artikel vor, so erhält man als Endergebnis (III)
\[
\begin{aligned}
t={}&\frac{rk\operatorname{tang}B\mathord{.}\sin T}{a}(1+\frac{5ee+(4ee-14e^4)\sin B^2+5e^4\sin B^4}{24k^4}\mathord{.}bb\\
&\qquad+\frac{k^4}{12aa\cos\varphi^2}\mathord{.}rr+\tfrac1{24}tt)
\end{aligned}
\]

\indent Die drei Formeln I, II, III enthalten im Wesentlichen die Lösung unserer Aufgabe. Aus ihrer Herleitung folgt unmittelbar, dass sie bis zur dritten Ordnung einschließlich genau sind. Dass ihre Genauigkeit tatsächlich noch eine Ordnung weiter reicht, der Fehler also bei jeder der Formeln von fünfter Ordnung ist, ließe sich leicht durch einige ergänzende Zwischenentwicklungen zeigen. Man könnte es auch daran zeigen, dass in den Ausdrücken ihrer Natur nach keine Größen gerader Ordnung vorkommen können. Darauf gehe ich hier jedoch nicht weiter ein, da die \textit{zweite} Herleitung der Formeln in den folgenden Artikeln (26---32) von selbst dasselbe Ergebnis liefert.

\Article{26}
Diese Untersuchung ist als eigenständig und unabhängig von allem Vorhergehenden zu betrachten. Der Übersicht halber und um Unklarheiten zu vermeiden, werden daher alle verwendeten Bezeichnungen bei ihrer Einführung erklärt. Die meisten Buchstaben aus der ersten Herleitung behalten ihre dortige Bedeutung. Zwei von ihnen (\(u\) und \(s\)), die dort nur Hilfsgrößen darstellen%
